\Qhead{1.6}{Автоколебательные системы с двумя и более степенями свободы}{Автоколебательные системы с двумя и более степенями свободы. Взаимная синхронизация колебаний двух генераторов.}
\Subfile{1.6.1}{sections/01_oscillations/q06/s01}
\Subfile{1.6.2}{sections/01_oscillations/q06/s02}

\recap{Несколько мод общего генератора конкурируют за энергию. Два самостоятельных генератора могут согласовать фазы: разность фаз подчиняется уравнению Адлера с суммой взаимных воздействий. Существование режима и его устойчивость проверяются раздельно.}
\board{Два связанных контура и общий активный элемент; само- и перекрёстное насыщение; $K_1,K_2$ и общая частота.}
\oral{Почему «слабая конкуренция» без проверки положительности интенсивностей не гарантирует двухмодовую генерацию?}
\bait{Можно оставить односторонний предел связи: если один генератор не испытывает обратного действия, взаимная синхронизация превращается во внешний захват.}
\sources{Балакший~\cite{balakshy}, PDF-с.~16 и 19, листы~14 и 17: связанные контуры, переключение мод и общая постановка. Вопрос~6 коллеги задаёт порядок объяснения; условия существования, устойчивости и определения $K_j$ уточнены самостоятельным расчётом.}
